\section{Spherical geometry and the stability inequality}
\label{geo:section}
\label{sec:geometry}

Let $\Omega_o\subset S^{m_o}$, $m_o\in\{4,5\}$, be the original
component supplied by Proposition~\ref{red:original-link} and
Lemma~\ref{red:nonflat-component}. Its closure is a compact smooth
manifold with nonempty boundary. The function $g_o$ is smooth up to that
boundary, is positive in $\Omega_o$, and satisfies
\[
 \Delta g_o=-m_o g_o,\qquad
 g_o=0,\quad \nabla g_o=N_o\quad\hbox{on }\partial\Omega_o,
\]
where $N_o$ is the inward unit normal. Moreover,
$\nabla^2g_o+g_o\Id$ does not vanish identically.
We give the identities in a general tangent dimension $m$. The compact
original component first provides strict gradient and boundary-sign
bounds; we then transport these bounds to the working link and derive
its stability inequality.
All spherical derivatives below are covariant, and $\Delta$ on tensors
is the rough Laplacian $\tr\nabla^2$. Repeated indices in an orthonormal
frame are summed.

\subsection{The spherical Hessian and its boundary values}

For a smooth solution of $\Delta g=-mg$ on a domain in $S^m$, set
\begin{equation}\label{geo:data}
 p=\nabla g,\qquad T=\nabla^2g+g\Id,\qquad
 B_{ijk}=\nabla_iT_{jk},\qquad w=|T|,\qquad s^2=g^2+|p|^2.
\end{equation}
The notation $s^2$ denotes the displayed smooth function, including at
points where it vanishes. If $u(ry)=rg(y)$ is the homogeneous extension,
then in the radial and tangential orthonormal frame
\begin{equation}\label{geo:ambient-hessian}
 D^2u(ry)=r^{-1}
 \begin{pmatrix}0&0\\0&T(y)\end{pmatrix}.
\end{equation}
Indeed, Euler's identity gives $D^2u\,y=0$, and differentiation of
$\nabla_{\R^{m+1}}u=g y+p$ along a tangent vector gives the tangential
block $\nabla^2g+g\Id$ at radius one. Homogeneity then gives the factor
$r^{-1}$.
The same gradient formula identifies $s^2$ with
$|\nabla_{\R^{m+1}}u|^2$ on the unit sphere.

\begin{proposition}[Interior identities]\label{geo:interior-identities}
The tensor $T$ is symmetric and trace-free, and $B$ is fully symmetric
and trace-free in every pair of slots. In particular $\diver T=0$.
The following identities hold:
\begin{align}
 \nabla_lB_{ijk}-\nabla_iB_{ljk}
 &=\delta_{lj}T_{ik}+\delta_{lk}T_{ij}
   -\delta_{ij}T_{lk}-\delta_{ik}T_{lj},
 \label{geo:commutator}\\
 \Delta T&=mT,
 &\Delta w^2&=2m w^2+2|B|^2,
 \label{geo:bochner}\\
 \nabla_l p_i&=T_{li}-g\delta_{li},
 &\nabla_l g&=p_l,
 \label{geo:first-jets}\\
 \nabla s^2&=2Tp,
 &\Delta s^2&=2w^2.
 \label{geo:s-identities}
\end{align}
On the open set $\{w>0\}$ one also has
\begin{equation}\label{geo:w-gradient}
 \nabla_l w=w^{-1}T_{jk}B_{ljk},\qquad |\nabla w|\le |B|.
\end{equation}
The inequality
$|\nabla(w^2)|^2\le4w^2|B|^2$ holds everywhere, including at $w=0$.
\end{proposition}

\begin{proof}
The trace equation is $\tr T=\Delta g+mg=0$.
On the unit sphere the commutation rule for the covector $p$ gives
\[
 \nabla_i(\nabla^2g)_{jk}-\nabla_j(\nabla^2g)_{ik}
 =\delta_{ik}p_j-\delta_{jk}p_i.
\]
Adding the derivatives of $g\Id$ cancels the right side. Thus $T$ is
Codazzi: $\nabla_iT_{jk}=\nabla_jT_{ik}$.
Together with symmetry in $j,k$, this proves full symmetry of $B$.
Differentiating $\tr T=0$ shows $B_{ijj}=0$; full symmetry gives all
other traces and hence $\diver T=0$.

For clarity, the same curvature convention on a covariant two-tensor is
\[
 ([\nabla_l,\nabla_i]T)_{jk}
 =\delta_{lj}T_{ik}+\delta_{lk}T_{ij}
  -\delta_{ij}T_{lk}-\delta_{ik}T_{lj}.
\]
This proves \eqref{geo:commutator}. Substitute $(l,i,j,k)=(i,j,i,k)$
there and sum over $i$. By symmetry and the zero traces of $B$, the
left side is $\Delta T_{jk}$, whereas the right side is
\[
 mT_{jk}+T_{jk}-T_{jk}-\delta_{jk}\tr T=mT_{jk}.
\]
The product rule for $|T|^2$ now proves \eqref{geo:bochner}.
The definitions give \eqref{geo:first-jets} and
\[
 \nabla_l s^2=2gp_l+2p_i(T_{li}-g\delta_{li})=2T_{li}p_i.
\]
Taking the divergence and using $\diver T=0$ and $\tr T=0$ gives
$\Delta s^2=2T_{li}(T_{li}-g\delta_{li})=2w^2$.
Finally, differentiating $w^2$ and applying Cauchy--Schwarz proves
both gradient estimates.
\end{proof}

\begin{lemma}[Strict gradient bound on the original component]
\label{geo:gradient-bound}
On $\Omega_o$ one has $s_o^2<1$. At every boundary point, put
\begin{equation}\label{geo:H-definition}
 H_o=-T_o(N_o,N_o).
\end{equation}
Then
\begin{equation}\label{geo:strict-boundary-sign}
 H_o>0,\qquad T_op_o=-H_op_o,\qquad w_o>0
 \quad\hbox{on }\partial\Omega_o.
\end{equation}
Consequently $w_o$ has a positive minimum on $\partial\Omega_o$,
and its zero set is a compact subset of $\Omega_o$.
\end{lemma}

\begin{proof}
The boundary conditions give $s_o^2=1$ on $\partial\Omega_o$.
By \eqref{geo:s-identities}, $s_o^2$ is subharmonic.
The maximum principle on the compact closure gives $s_o^2\le1$.
If equality held at an interior point, the strong maximum principle
would make $s_o^2$ constant on the connected component, and then
$\Delta s_o^2=2w_o^2$ would imply $T_o\equiv0$. Thus $s_o^2<1$ inside.

The smooth boundary has the interior sphere property at each point.
The Hopf boundary lemma therefore gives
$\partial_{N_o}s_o^2<0$. Since $p_o=N_o$ there,
\eqref{geo:s-identities} identifies this derivative as
$2T_o(N_o,N_o)$, proving $H_o>0$.
For a tangent vector $\tau$ to $\partial\Omega_o$, tangential
differentiation of $s_o^2=1$ similarly gives $T_o(\tau,N_o)=0$.
This proves the eigenvector assertion and $w_o\ge H_o>0$.
Continuity and compactness of the boundary give a positive minimum
of $w_o$ and a neighborhood of the boundary disjoint from its zero set.
\end{proof}

The remaining boundary formulas are local. Thus they apply on the
original boundary and, once the strict sign has been transported,
on the regular side of the suspension.

\begin{proposition}[Boundary jets]\label{geo:boundary-jets}
Suppose $g=0$, $p=N$, and $T(N,N)=-H<0$ on a smooth boundary portion,
where $N$ is inward. Let Greek indices be tangent to this boundary and
put
\[
 \mathfrak h_{\alpha\beta}
 =\langle\nabla_{\alpha}N,e_\beta\rangle.
\]
Then $T_{N\alpha}=0$, $T_{\alpha\beta}=\mathfrak h_{\alpha\beta}$,
and $H=\tr\mathfrak h$; in particular $H$ is a sum of principal
curvatures with the specified convention, without division by $m-1$.
If $\nabla^{\partial}$ is the boundary connection, all components of
$B$ are described, up to permutation of slots, by
\begin{align}
 B_{\alpha\beta\gamma}
   &=(\nabla^{\partial}_{\alpha}\mathfrak h)_{\beta\gamma},
 &B_{NN\alpha}&=-\partial_\alpha H,
 \label{geo:even-boundary-jets}\\
 B_{N\alpha\beta}
   &=T_{NN}T_{\alpha\beta}-(T^2)_{\alpha\beta},
 &B_{NNN}&=w^2.
 \label{geo:odd-boundary-jets}
\end{align}
The components with an odd number of normal slots are thus determined
by the pointwise tensor $T$; the other components involve tangential
derivatives of curvature. The tensor $V$ below encodes this odd-normal
part after normalization.
Here $w>0$. Define the normalized tensors and eigenvalue
\[
 A=T/w,\qquad \lambda=T_{NN}/w=-H/w,
 \qquad S_0=\lambda A-A^2,
\]
and the symmetric tensor
\begin{equation}\label{geo:boundary-V}
 V_{ijk}=p_i(S_0)_{jk}+p_j(S_0)_{ki}+p_k(S_0)_{ij}
           +p_ip_jp_k.
\end{equation}
Then $|p|=1$, $Ap=\lambda p$, $\lambda<0$, $S_0p=0$,
$\tr S_0=-1$, and $V$ is trace-free. The tensors $B/w^2$ and $V$
agree in every component with an odd number of normal slots. Moreover,
\begin{equation}\label{geo:normal-w}
 V_{Njk}=(S_0+p\otimes p)_{jk},\qquad
 \partial_N w=w^2\bigl(2\lambda-\tr A^3\bigr).
\end{equation}
\end{proposition}

\begin{proof}
On the boundary, tangential differentiation of $|p|^2=1$ gives
$T_{N\alpha}=0$, and $\nabla_\alpha p=T e_\alpha$ because $g=0$.
Since $p=N$ along the boundary, this proves
$T_{\alpha\beta}=\mathfrak h_{\alpha\beta}$.
The trace equation for $T$ gives $\tr\mathfrak h=-T_{NN}=H$.
In differentiating three tangential entries of $T$, the normal
connection terms pair with $T(N,e_\alpha)=0$. This yields the first
formula in \eqref{geo:even-boundary-jets}.
Likewise,
\[
 \partial_\alpha T(N,N)
 =B_{\alpha NN}+2T(\nabla_\alpha N,N)=B_{\alpha NN},
\]
which proves the second formula by symmetry of $B$.

For a tangent field $e_\alpha$, differentiate $T(N,e_\alpha)=0$
along $e_\beta$. Since
$\langle N,\nabla_\beta e_\alpha\rangle
=-\mathfrak h_{\beta\alpha}$, this gives
\[
 0=B_{\beta N\alpha}+(T^2)_{\beta\alpha}
     -T_{NN}T_{\beta\alpha}.
\]
This proves the first formula in \eqref{geo:odd-boundary-jets}.
Taking its tangential trace gives
\[
 \sum_\alpha B_{N\alpha\alpha}
 =T_{NN}\sum_\alpha T_{\alpha\alpha}
    -\sum_{\alpha,\beta}T_{\alpha\beta}^2
 =-T_{NN}^2-\sum_{\alpha,\beta}T_{\alpha\beta}^2=-w^2.
\]
The zero trace of $B$ then gives $B_{NNN}=w^2$.

Normalization gives $S_0p=0$ and $\tr S_0=-|A|^2=-1$.
Consequently $V_{ijj}=p_i\tr S_0+2(S_0p)_i+p_i=0$.
In an orthonormal frame with first vector $p=N$, one has
$V_{N\alpha\beta}=(S_0)_{\alpha\beta}$ and $V_{NNN}=1$.
These are exactly the normalized identities
\eqref{geo:odd-boundary-jets}, and $V_{Njk}=S_{0,jk}+p_jp_k$.
Finally, $T$ has only its normal-normal and tangential-tangential
blocks, so the odd-normal formulas account for the entire contraction
$T_{jk}B_{Njk}$. Thus
\[
 \partial_Nw=w^{-1}T_{jk}B_{Njk}
 =w^2 A_{jk}(S_0+p\otimes p)_{jk}
 =w^2(2\lambda-\tr A^3),
\]
as asserted.
\end{proof}

\subsection{Transport to a suspension}

\begin{proposition}[Suspended data]\label{geo:suspension-data}
If $m_o=4$, write $\rho=\sin t$, $c=\cos t$ and parameterize the
regular part of the suspended component by
\[
 (t,y)\longmapsto (\rho y,c),\qquad
 0<t<\pi,\quad y\in\Omega_o.
\]
Its metric is $dt^2+\rho^2 ds_{S^4}^2$, and $\widetilde g=\rho g_o$.
With $E_0=\partial_t$ and $E_\alpha=\rho^{-1}e_\alpha$ for a lifted
orthonormal frame on $S^4$, the data satisfy
\begin{align}
 \widetilde p_0&=c g_o,
 &\widetilde p_\alpha&=(p_o)_\alpha,
 &\widetilde s^2&=s_o^2,
 \label{geo:suspended-first-jets}\\
 \widetilde T_{00}&=0,
 &\widetilde T_{0\alpha}&=0,
 &\widetilde T_{\alpha\beta}&=\rho^{-1}(T_o)_{\alpha\beta}.
 \label{geo:suspended-T}
\end{align}
The only possibly nonzero components of $\widetilde B$ are
\begin{equation}\label{geo:suspended-B}
 \widetilde B_{\alpha\beta\gamma}
   =\rho^{-2}(B_o)_{\alpha\beta\gamma},\qquad
 \widetilde B_{0\alpha\beta}
 =\widetilde B_{\alpha0\beta}
 =\widetilde B_{\alpha\beta0}
   =-c\rho^{-2}(T_o)_{\alpha\beta}.
\end{equation}
In particular,
\begin{equation}\label{geo:suspended-norms}
 \widetilde w=\rho^{-1}w_o,\qquad
 |\widetilde B|^2
 =\rho^{-4}\bigl(|B_o|^2+3c^2w_o^2\bigr),\qquad
 \widetilde H=\rho^{-1}H_o
 \quad\hbox{on the regular side.}
\end{equation}
The volume and side-boundary measures are
\begin{equation}\label{geo:suspended-measures}
 dV=\rho^4\,dt\,dV_o,\qquad
 d\sigma=\rho^3\,dt\,d\sigma_o.
\end{equation}
Every local identity of Proposition~\ref{geo:interior-identities}
holds with $m=5$, and Proposition~\ref{geo:boundary-jets} holds on
the regular side. Also $\widetilde s^2<1$ inside and
$\widetilde H,\widetilde w>0$ on that side.
\end{proposition}

\begin{proof}
At a point where the base frame is normal, the warped metric has
connection
\[
 \nabla_{E_0}E_0=\nabla_{E_0}E_\alpha=0,\qquad
 \nabla_{E_\alpha}E_0=(c/\rho)E_\alpha,\qquad
 \nabla_{E_\alpha}E_\beta=-(c/\rho)\delta_{\alpha\beta}E_0.
\]
The horizontal part of the last expression at other points is the
base connection with factor $\rho^{-1}$.
Differentiating $\widetilde g=\rho g_o$ gives
\eqref{geo:suspended-first-jets}. Its Hessian has entries
\[
 (\nabla^2\widetilde g)_{00}=-\rho g_o,
 \quad (\nabla^2\widetilde g)_{0\alpha}=0,
 \quad (\nabla^2\widetilde g)_{\alpha\beta}
 =\rho^{-1}(\nabla_o^2 g_o)_{\alpha\beta}
   +(c^2/\rho)g_o\delta_{\alpha\beta}.
\]
Adding $\widetilde g\Id$ proves \eqref{geo:suspended-T}.
Differentiating that tensor with the displayed connection gives
\eqref{geo:suspended-B}; entries with two or three zero slots vanish.
Each term with exactly one zero slot occurs in three positions,
which proves the factor $3$ in \eqref{geo:suspended-norms}.

The inward normal to the side is the horizontal lift of $N_o$.
Thus the normal-normal entry of \eqref{geo:suspended-T} gives
$\widetilde H=\rho^{-1}H_o$.
The metric on the side is $dt^2+\rho^2 ds_{\partial\Omega_o}^2$,
proving \eqref{geo:suspended-measures}.
The trace of \eqref{geo:suspended-T} is zero, so
$\Delta\widetilde g=-5\widetilde g$.
The interior identities follow locally from
Proposition~\ref{geo:interior-identities}.
The bounds for $s^2$, $H$ and $w$ follow from the original compact
component and \eqref{geo:suspended-first-jets}--\eqref{geo:suspended-norms};
no maximum principle at a tip is used. The local boundary hypotheses
and hence Proposition~\ref{geo:boundary-jets} hold on the regular side.
\end{proof}

We now use the unadorned notation $\Omega,g,p,T,B,w,H$ for the
working component in $S^5$ defined at the end of
Section~\ref{sec:reductions}. It is either
$\Omega_o$ when $m_o=5$, or the above suspension when $m_o=4$.
All boundary integrals mean integrals over the regular boundary, with
the two tips omitted in the suspended case. The original data retain
the subscript $o$ when needed.

\subsection{Second variation under compact perturbations}

We derive stability directly from the minimizing property, including
its sign and test class. The argument works in either original
ambient dimension and in the six-dimensional suspension.

\begin{proposition}[Euclidean stability on a component]
\label{geo:euclidean-stability}
Let $D$ be the selected positivity component of a global minimizing
cone in $\R^{m+1}$. Let $\phi$ be any real smooth function up to the
regular boundary of $D$, with compact support in that regular closure
away from the vertex and from any other singular points. Then
\begin{equation}\label{geo:euclidean-inequality}
 \int_D |\nabla\phi|^2\,dx
 \ge \int_{\partial D}H_E\phi^2\,d\sigma_E,
 \qquad H_E=-D^2u(N_E,N_E),
\end{equation}
where $N_E$ is the inward Euclidean unit normal. No restriction is
imposed on the boundary trace or normal derivative of $\phi$.
\end{proposition}

\begin{proof}
Every regular free-boundary point has a neighborhood where the
positive phase is the single side belonging to $D$. Compactness of
the boundary portion meeting $\supp\phi$ gives finitely many such
neighborhoods and a collar separated from every other positivity
component and every singular point. Choose the collar over a slightly
larger boundary portion, so that the extension of $\phi$ can vanish
near its lateral edge. Extend $\phi$ smoothly across this boundary
with compact support in these neighborhoods. Also extend the positive
side of $u$ smoothly to a function $F$ on the collar, with its interior
jets unchanged. Since $u=0$ and $\nabla u=N_E$ on the boundary,
shrinking the collar makes $F$ negative on its exterior side and
positive on its interior side. Such an extension need only be smooth
to finite order sufficient for the expansions below; smooth boundary
regularity supplies this, and a partition of unity joins the local
extensions while preserving the interior function and its jets.

For a real parameter $\varepsilon$ of either sufficiently small sign,
use $(F+\varepsilon\phi)_+$ in the collar,
$u+\varepsilon\phi$ on the rest of the affected portion of $D$, and
$u$ elsewhere. These definitions agree: $F=u$ on the interior of the
collar, and on the compact interior support outside a smaller collar
$u$ has a positive minimum. Near the exterior and lateral edges of
the modification the function agrees with $u$.
The resulting competitor $u_\varepsilon$ is nonnegative, belongs to
$H^1_{\mathrm{loc}}$, and has compactly supported difference from $u$ in a containing
ball. It is therefore admissible for the original global minimizer.
This construction uses no separate minimizing assertion for $D$.

We compute the change in energy in such a ball. First integrate the
perturbed density over the fixed original domain $D$, and then account
for the signed boundary layer where its positivity set changes.
The fixed-domain contribution is
\begin{align}
 &2\varepsilon\int_D\langle\nabla u,\nabla\phi\rangle\,dx
       +\varepsilon^2\int_D|\nabla\phi|^2\,dx\notag\\
 &\hspace{15mm}=-2\varepsilon\int_{\partial D}\phi\,d\sigma_E
       +\varepsilon^2\int_D|\nabla\phi|^2\,dx.
 \label{geo:fixed-domain-variation}
\end{align}
Here harmonicity eliminates the interior term and the outward normal
is $-N_E$, so $\partial_{-N_E}u=-1$.
All integrations by parts have compact support on the regular portion.

Use coordinates $(y,\xi)$ obtained by flowing from the original
boundary under
\[
 V_F=\frac{\nabla F}{|\nabla F|^2},\qquad \xi=F.
\]
The denominator is nonzero on a sufficiently small compact collar.
At $\xi=0$, $V_F=N_E$. If $j(y,\xi)$ is the volume Jacobian relative
to $d\sigma_E(y)\,d\xi$, then
\begin{equation}\label{geo:flow-jacobian}
 j(y,0)=1,\qquad
 \partial_\xi j(y,0)=\diver V_F(y,0)
 =\Delta F-2F_{N_EN_E}=2H_E.
\end{equation}
The equality $\Delta F=0$ on the original boundary follows from
matching the interior second derivatives; harmonicity outside $D$
is unnecessary. Moreover,
\[
 \left.\partial_\xi|\nabla F|^2\right|_{\xi=0}
 =2F_{N_EN_E}=-2H_E,
\]
and hence
\begin{equation}\label{geo:edge-density-derivative}
 \left.\partial_\xi\bigl[(1+|\nabla F|^2)j\bigr]
 \right|_{\xi=0}=2H_E.
\end{equation}

Write $\phi_0=\phi(y,0)$ and
$\phi_N=\partial_{N_E}\phi(y,0)$.
The new zero surface has coordinate $b_\varepsilon(y)$ determined by
$b_\varepsilon+\varepsilon\phi(y,b_\varepsilon)=0$.
The implicit function theorem and
$\partial_\xi\phi(y,0)=\phi_N$ give, uniformly on the compact patch,
\begin{equation}\label{geo:moving-edge}
 b_\varepsilon=-\varepsilon\phi_0
      +\varepsilon^2\phi_0\phi_N+O(|\varepsilon|^3).
\end{equation}
Throughout the layer $|\xi|\le C|\varepsilon|$, its perturbed energy
density is
\begin{equation}\label{geo:layer-density}
 (1+|\nabla(F+\varepsilon\phi)|^2)j
 =2+2\varepsilon\phi_N+2H_E\xi
      +O(\varepsilon^2+|\varepsilon\xi|+\xi^2).
\end{equation}
Integrating this density from $b_\varepsilon$ to $0$ is the signed
correction, whether the phase expands or contracts. The terms of
order at most two are
\begin{align}
 \int_{b_\varepsilon}^{0}
       (1+|\nabla(F+\varepsilon\phi)|^2)j\,d\xi
 &=-2b_\varepsilon-2\varepsilon\phi_N b_\varepsilon
      -H_E b_\varepsilon^2+O(|\varepsilon|^3)\notag\\
 &=2\varepsilon\phi_0
   +\varepsilon^2\bigl(-2\phi_0\phi_N
          +2\phi_0\phi_N-H_E\phi_0^2\bigr)
       +O(|\varepsilon|^3)\notag\\
 &=2\varepsilon\phi_0-\varepsilon^2H_E\phi_0^2
       +O(|\varepsilon|^3).
 \label{geo:layer-cancellation}
\end{align}
This displays the cancellation of the normal-derivative terms.
The uniform remainder is integrable over the compact affected
boundary. Adding \eqref{geo:fixed-domain-variation} also cancels the
linear terms, leaving
\[
 J(u_\varepsilon)-J(u)
 =\varepsilon^2\left(
      \int_D|\nabla\phi|^2\,dx
        -\int_{\partial D}H_E\phi^2\,d\sigma_E\right)
       +o(\varepsilon^2).
\]
The energies here are in a ball containing the modification.
Minimality and division by $\varepsilon^2$ prove
\eqref{geo:euclidean-inequality}.
\end{proof}

The angular constant comes from the sharp radial Hardy threshold,
which also enters the cone instability criteria in
\cite[Propositions~2.1--2.2]{JerisonSavin2015}. We include the radial
calculation to specify the compact supports and the order of limits.

\begin{corollary}[Spherical stability]\label{geo:spherical-stability}
For every real function $v$ smooth up to the regular boundary of the
working component $\Omega\subset S^5$, with compact support in its
closure away from the tips if present,
\begin{equation}\label{geo:spherical-inequality}
 \int_{\partial\Omega}H v^2\,d\sigma
 \le \int_\Omega\bigl(|\nabla v|^2+4v^2\bigr)\,dV.
\end{equation}
The trace of $v$ on the regular boundary is arbitrary.
\end{corollary}

\begin{proof}
More generally, consider a link component $\Omega\subset S^m$ with
angular tests supported in its regular closure, and write
$D=\{ry:r>0,\ y\in\Omega\}$.
On its regular side the normal is tangent to $S^m$, and
\eqref{geo:ambient-hessian} gives $H_E(ry)=r^{-1}H(y)$.
Polar volume has factor $r^m\,dr$ and side area has factor
$r^{m-1}\,dr$.
For $\zeta\in C_c^\infty((0,\infty))$ and an allowed angular test $v$,
put $\phi(ry)=\zeta(r)v(y)$ in
Proposition~\ref{geo:euclidean-stability}. It gives
\begin{align*}
 \left(\int_0^\infty r^{m-2}\zeta^2\,dr\right)
       \int_{\partial\Omega}Hv^2\,d\sigma
 &\le
 \left(\int_0^\infty r^{m-2}\zeta^2\,dr\right)
       \int_\Omega|\nabla v|^2\,dV\\
 &\quad+
 \left(\int_0^\infty r^m|\zeta'|^2\,dr\right)
       \int_\Omega v^2\,dV.
\end{align*}
For each fixed $\zeta$ and $v$, the Euclidean support lies in a finite
annulus and is separated from the singular set and other components,
so the preceding second variation applies.

Set $\beta=(m-1)/2$. For $f_\zeta(r)=r^\beta\zeta(r)$, integration of
the cross term gives the radial Hardy identity
\[
 \int_0^\infty r^m|\zeta'|^2\,dr
  -\beta^2\int_0^\infty r^{m-2}\zeta^2\,dr
 =\int_0^\infty r|f_\zeta'(r)|^2\,dr\ge0.
\]
To approach its constant, choose nonzero
$\psi\in C_c^\infty(\R)$ and set
$\zeta_k(r)=r^{-\beta}\psi((\log r)/k)$.
The substitution $\tau=(\log r)/k$ yields
\begin{equation}\label{geo:hardy-limit}
 \frac{\int_0^\infty r^m|\zeta_k'|^2\,dr}
      {\int_0^\infty r^{m-2}\zeta_k^2\,dr}
 =\beta^2+\frac1{k^2}
          \frac{\int_\R|\psi'|^2\,d\tau}
               {\int_\R\psi^2\,d\tau}.
\end{equation}
Indeed the numerator before division is
$k\int_\R(-\beta\psi+k^{-1}\psi')^2\,d\tau$, and
$\int\psi\psi'=0$ by compact support.
Divide the separated stability inequality by its positive radial
factor and let $k\to\infty$. In tangent dimension $m=5$, one has
$\beta^2=4$, proving \eqref{geo:spherical-inequality}.
The same computation on an original link with $m_o=4$ would give
$9/4$; the working link always has $m=5$.

The order of limits is fixed: for each angular test and each finite
$k$, first take $\varepsilon\to0$ in the second variation, and then
take $k\to\infty$. Angular cutoffs used later are removed only after
\eqref{geo:spherical-inequality} has been obtained for each fixed
smooth test. The collar and the allowable size of $\varepsilon$ may
depend on these fixed supports.
\end{proof}
